\section{Введение: ширина как ресурс}\label{sec:intro}

\subsection{Постановка и критерий}

Через $\chi\bigl(\R^n,[1,\ell]\bigr)$ обозначим наименьшее число цветов
раскраски пространства $\R^n$, при которой никакие две одноцветные точки не
находятся на расстоянии из отрезка $[1,\ell]$, $\ell\ge1$. При $\ell=1$ это
классическое хроматическое число $\chi(\R^n)$ (проблема Нельсона--Хадвигера;
о состоянии вопроса см.~\cite{Soifer,deGrey}), интервальная версия
систематически изучается с~\cite{Ivanov2011}; в плоскости она же
рассматривалась в форме запрета $[1-\varepsilon,1+\varepsilon]$
\cite{Exoo2005,ChJW,deGreyParts,Parts2023}. Запрет более широкого множества
расстояний требует не меньше цветов, поэтому
$\chi(\R^n)\le\chi\bigl(\R^n,[1,\ell]\bigr)$ при всех $\ell\ge1$.

Все известные верхние оценки $\chi(\R^n)$ в малых размерностях получены
периодическими раскрасками по разбиению Вороного решётки: для решётки
$\Lambda\subset\R^n$ и подрешётки $\Gamma\subseteq\Lambda$ индекса
$k=[\Lambda:\Gamma]$ ячейка Вороного $V_\Lambda(v)$ окрашивается цветом
класса $v+\Gamma$. Пусть $V_0=V_\Lambda(0)$,
\[
D(v)=\dist\bigl(V_0,\,v+V_0\bigr),\qquad
D(\Gamma)=\min_{v\in\Gamma\setminus\{0\}}D(v),\qquad
d(\Lambda,\Gamma)=\frac{D(\Gamma)}{\diam V_0}.
\]
Одноцветные точки лежат либо в одной ячейке (расстояние $\le\diam V_0$),
либо в ячейках с центрами, отличающимися на $v\in\Gamma\setminus\{0\}$
(расстояние $\ge D(\Gamma)$), так что промежуток
$(\diam V_0,D(\Gamma))$ свободен.

\begin{proposition}[критерий; {\cite[предложение~1]{Bounds}}]\label{prop:crit}
Если $d(\Lambda,\Gamma)>1$, то $\chi\bigl(\R^n,[1,\ell]\bigr)\le[\Lambda:\Gamma]$
при всех $\ell<d(\Lambda,\Gamma)$; в частности $\chi(\R^n)\le[\Lambda:\Gamma]$.
\end{proposition}

Величину $d$ будем называть \emph{нормированной шириной} запрещённого
интервала, или просто \emph{шириной} конструкции. Вычислять $D(v)$ помогает
центральная симметрия ячейки Вороного и всех её фасет
(лемма Вороного~\cite{Voronoi}, далее --- лемма~\ref{lem:sym}):

\begin{lemma}[о середине; {\cite[лемма~2]{Ivanov2011}}]\label{lem:mid}
$D(v)=2\,\dist\bigl(v/2,\,V_0\bigr)$.
\end{lemma}

\begin{lemma}[Вороной]\label{lem:sym}
Ячейка $V_0$ --- выпуклый центрально-симметричный многогранник, все её фасеты
центрально-симметричны, и $\diam V_0=2R$, где $R$ --- радиус покрытия.
\end{lemma}

Минимум $D(\Gamma)$ берётся по конечному окну: так как
$D(v)\ge\|v\|-\diam V_0$, улучшить уже найденное значение $D_*$ могут лишь
векторы с $\|v\|<D_*+\diam V_0$.

В статье~\cite{Bounds} ширина играла роль отчёта о качестве конструкции: там
доказаны пять новых оценок
\[
\chi(\R^4)\le43,\quad\chi(\R^5)\le132,\quad\chi(\R^7)\le1029,\quad
\chi(\R^9)\le9604,\quad\chi(\R^{10})\le45619,
\]
и у каждой указана ширина запрещённого интервала. Ключом к двум последним
служит \emph{продуктовое правило}, которое мы будем использовать и здесь:

\begin{proposition}[продуктовое исчисление; {\cite[предложение~4]{Bounds}}]
\label{prop:product}
Пусть $\Lambda=\bigoplus_{i=1}^{s}\Lambda_i$ --- ортогональная прямая сумма
решёток $\Lambda_i\subset\R^{n_i}$, $\Gamma=\bigoplus_i\Gamma_i$,
$d_i=d(\Lambda_i,\Gamma_i)$, $k_i=[\Lambda_i:\Gamma_i]$. Тогда после
независимого масштабирования блоков раскраска по $\Gamma$ допустима для
интервала $[1,\ell]$ с $\ell=\bigl(\sum_i 1/d_i^2\bigr)^{-1/2}$, то есть
$\chi(\R^n,[1,\ell])\le\prod_i k_i$ при $n=\sum_i n_i$, как только
$\sum_i 1/d_i^2\le1$; оптимум достигается при $\diam^2V_0(\Lambda_i)\propto1/d_i^2$.
\end{proposition}

Правило превращает ширину из отчётной характеристики в \emph{расходуемый
ресурс}: величина $1/d_i^2$ --- цена блока, а бюджет равен единице. Настоящая
статья посвящена ширине как самостоятельному предмету: сколько её даёт
конструкция, как она расходуется при произведении и ламинировании и где лежат
строгие либо вычислительные пределы метода.

\subsection{Результаты}

Точные результаты статьи собраны в таблице~\ref{tab:main}; в ней и далее
\stP{} --- теорема, \stC{} --- точный рациональный сертификат (утверждение
сведено к конечному списку неравенств между явно выписанными дробями,
проверенных в точной арифметике; плавающая точка могла участвовать в выборе
кандидата, но не в проверке), \stN{} --- \emph{численный кандидат}: величина
получена вычислениями с плавающей точкой и доказательной силы не имеет.
Численные кандидаты в статье всюду названы кандидатами и в утверждения не
входят.

\begin{table}[htbp]\centering\small
\caption{Точные результаты статьи. В столбце $\ell$ стоит доказанная ширина
запрещённого интервала $[1,\ell]$; запись $<x$ означает, что оценка доказана
при каждом $\ell<x$.}
\label{tab:main}
\begin{tabular}{@{}clllc@{}}
\toprule
$n$ & цветов & $\ell$ & утверждение & \\
\midrule
$2$ & $8$ & $1{,}441350$ & мозаичная раскраска, \S\ref{ssec:quant} & \stC\\
$2$ & $15$ & $2{,}281534$ & мозаичная раскраска, \S\ref{ssec:quant} & \stC\\
$3$ & $15$ & $102659/100000$ & семейство $G(\alpha)$, \S\ref{sec:stair} & \stC\\
$3$ & $\ge15$ & --- & неулучшаемость $15$ на ОЦК, теорема~\ref{thm:a3opt} & \stP\\
$4$ & $45$ & $1{,}015$ & решётка общего положения, теорема~\ref{thm:main} & \stC\\
$4$ & $48$ & $1{,}0396$ & решётка общего положения, теорема~\ref{thm:main} & \stC\\
$6$ & $343$ & $<\sqrt{7/6}$ (точно) & тождество, теорема~\ref{thm:eis} & \stP\\
$8$ & $2401$ & $<\sqrt{7/6}$ (точно) & тождество, теорема~\ref{thm:eis} & \stP\\
$8$ & $\ge2401$ & --- & неулучшаемость на $E_8$, теорема~\ref{thm:e8opt} & \stP\\
$12$ & $3^{12}$ & $<\sqrt{3/2}$ & правило $d=2/\rho$ на $K_{12}$, \S\ref{ssec:k12} & \stP\\
$24$ & $7^{12}$ & $<\sqrt{7/6}$ (точно) & тождество, теорема~\ref{thm:eis} & \stP\\
$25$ & $4\cdot7^{12}$ & $<\sqrt{63/61}$ & произведение с $\Lambda_{24}$, \S\ref{ssec:dp} & \stP\\
$26$ & $19\cdot7^{12}$ & $<\sqrt{217/214}$ & произведение с $\Lambda_{24}$, \S\ref{ssec:dp} & \stP\\
любое & $\ge2^n$ & --- & пол мозаичного принципа, теорема~\ref{thm:floor} & \stP\\
\bottomrule
\end{tabular}
\end{table}

Содержание по разделам.

\begin{itemize}[leftmargin=1.4em,itemsep=2pt]
\item \S\ref{sec:identity}. \emph{Эйзенштейново тождество}
$D((3+\omega)\Lambda)=\sqrt{7/3}\,\lambda_1(\Lambda)$ доказано для всех
эйзенштейновых решёток (теорема~\ref{thm:eis}); отсюда точные ширины всех
конструкций $7^{n/2}$, в том числе $\sqrt{7/6}$ для $\chi(\R^{24})\le7^{12}$,
и $\alpha$-лестница для четырёх порядков ($\Z$, $\Z[i]$, $\Z[\omega]$,
кватернионы Гурвица) в форме равенства (следствие~\ref{cor:Uexact}). Правило
вещественного множителя $D(m\Lambda)=(m-1)\lambda_1$ даёт ширину
$\sqrt{3/2}$ для $\chi(\R^{12})\le3^{12}$ на решётке Коксетера--Тодда.
\item \S\ref{sec:ladders}. \emph{Лестницы ширин} в $\R^2$--$\R^6$ с точно
сертифицированными ступенями:
однопараметрическое семейство $G(\alpha)$ в $\R^3$, содержащее обе
$15$-раскраски Кулсона~\cite{Coulson} и достигающее максимума в корне
кубики, с сертификатом $\ell\le102659/100000$; компромисс
<<цвета\,$\leftrightarrow$\,ширина>> в $\R^4$ ($45$, $46$, $48$ цветов).
\item \S\ref{sec:lam}. \emph{Ламинирование} --- подъём раскраски $\R^{n-1}$
слоями --- и \emph{теорема о бюджете слоя}: избыток ширины базы есть в
точности бюджет радиуса покрытия на дополнительные измерения. Кусочный
сертификат диаметра сводит верхнюю границу радиуса покрытия слоёной решётки к
максимуму выпуклых квадратик по вершинам явных политопов; на нём стоят
численные кандидаты индексов $7203$ в $\R^9$ и $28812$ в $\R^{10}$, точного
сертификата у которых нет.
\item \S\ref{sec:floors}. \emph{Строгие полы индекса}: объёмный экран
Минковского, инрадиусный экран и \emph{оболочечный} экран, замечающий, что
норма вектора решётки пробегает не континуум, а оболочки. На двух классических
родителях пол совпадает с рекордом: $2401$ неулучшаемо на $E_8$
(теорема~\ref{thm:e8opt}) и $15$ --- на ОЦК (теорема~\ref{thm:a3opt}).
\item \S\ref{sec:tilings}. \emph{За пределами решёток}: раскраски
периодическими степенными (лагерровыми) мозаиками, пол $2^n$ для всего
класса (теорема~\ref{thm:floor}), структурная причина, по которой
нерешёточная свобода не окупается в чемпионских индексах, и точно
сертифицированные мозаичные ступени в плоскости.
\item \S\ref{sec:open}. Четыре открытых вопроса, которые эти результаты
ставят точнее прежнего.
\end{itemize}

Поисковые алгоритмы, полные лестницы (включая доминируемые ступени),
численные фронтиры и отрицательные экраны, а также большие матрицы и
протоколы всех проверок вынесены в полную версию работы~\cite{Full};
все точные сертификаты и код открыты в репозитории проекта.
